\documentclass[11pt]{article}
% =========================================================
%  學生講義 共用前言（以 XeLaTeX 編譯）
%  需要套件：xeCJK, tcolorbox（其餘多在 BasicTeX 內）
% =========================================================
\usepackage[a4paper,margin=2.3cm]{geometry}
\usepackage{amsmath,amssymb}
\usepackage{xcolor}
\usepackage{graphicx}

% ---- 中文字型（macOS 系統字型）----
\usepackage{xeCJK}
\setCJKmainfont{Songti TC}      % 內文：宋體
\setCJKsansfont{Heiti TC}       % 標題/強調：黑體
\xeCJKsetup{AutoFakeBold=2.2, AutoFakeSlant=0.2}

% ---- 版面 ----
\setlength{\parindent}{0pt}
\setlength{\parskip}{0.55em}
\linespread{1.15}
\renewcommand{\baselinestretch}{1.15}

% ---- 顏色 ----
\definecolor{cBlue}{HTML}{1f6feb}
\definecolor{cGray}{HTML}{6b7280}
\definecolor{cGrayBg}{HTML}{f3f4f6}
\definecolor{cPurp}{HTML}{7a3ff2}
\definecolor{cPurpBg}{HTML}{f5f1ff}

% ---- 標註框 ----
\usepackage[most]{tcolorbox}

% 就地補充新概念（灰底）：\begin{補充框}{標題}
\newtcolorbox{補充框}[1]{
  enhanced, breakable, arc=2pt, boxrule=0.6pt,
  colback=cGrayBg, colframe=cGray,
  left=9pt,right=9pt,top=7pt,bottom=7pt,
  before upper={\textbf{\sffamily #1}\par\vspace{3pt}},
}

% 延伸 / 開放問題（紫色左邊條）：\begin{延伸框}{標題}
\newtcolorbox{延伸框}[1]{
  enhanced, breakable, arc=2pt, boxrule=0pt,
  colback=cPurpBg, colframe=cPurpBg,
  borderline west={3pt}{0pt}{cPurp},
  left=10pt,right=9pt,top=7pt,bottom=7pt,
  before upper={\textbf{\sffamily 延伸閱讀｜#1}\par\vspace{3pt}},
}

% ---- 圖：置中、底下小字說明 ----
% \fig{檔名}{寬度}{說明}
\newcommand{\fig}[3]{%
  \begin{center}
  \includegraphics[width=#2\linewidth]{#1}\\[2pt]
  {\small\color{cGray} #3}
  \end{center}}

% ---- 章節標題用黑體 ----
\newcommand{\h}[1]{\sffamily #1}


\begin{document}

\begin{center}
{\sffamily\Huge\bfseries 選物I-5 週期運動}\\[3pt]
{\sffamily\large 普通高中 選修物理I（力學一）・學生講義}
\end{center}
\vspace{0.6em}

自然界存在大量會一再重複的運動：月球繞地球、彈簧上下振盪、鞦韆來回擺動、音叉振動。本章用\textbf{牛頓第二定律}定量描述這類\textbf{週期運動}，主要對象有兩個：\textbf{等速圓周運動（uniform circular motion）}與\textbf{簡諧運動（simple harmonic motion, SHM）}。最後會說明兩者是同一件事的兩個面向。

圓周運動是下一章萬有引力（行星、衛星繞行）的數學工具；簡諧運動則是日後波動、聲音、交流電以及量子力學中各種振盪行為的原型。本章同時用到向量、牛頓第二定律與三角函數三項先前所學的內容。

\medskip
本章主線：

\begin{center}
等速圓周運動 $\rightarrow$ 向心加速度與向心力 $\rightarrow$ 簡諧運動 $\rightarrow$ 簡諧與圓周的關係 $\rightarrow$ 單擺
\end{center}

\section*{\sffamily 5-1\quad 等速圓周運動}

\subsection*{\sffamily A.\ 為什麼「等速率」仍稱為「有加速度」}

\textbf{等速圓周運動}指物體沿圓形軌道運動，且\textbf{速率（speed，速度的大小）}保持不變。此處用「等速\textit{率}」而非「等速\textit{度}」，是因為\textbf{速度是向量}：圓周上每一點的速度都沿該點的\textbf{切線方向}，方向持續改變。

\begin{補充框}{速率不變，為什麼還有加速度？}
加速度量的不是「速率有沒有變快」，而是「\textbf{速度向量有沒有變}」。速度是向量，包含\textbf{大小（速率）與方向}兩部分；只要其中任何一個改變，速度就改變，就有加速度。等速圓周運動屬於「大小不變、方向持續改變」，因此仍有加速度。

加速度的定義是
$$\vec a=\frac{\Delta\vec v}{\Delta t}.$$
分子是 $\Delta\vec v$（速度\textbf{向量}的變化），而非 $\Delta v$（速率這個\textbf{數值}的變化）。把圓周上相鄰兩點的速度 $\vec v_1$、$\vec v_2$ 平移到同一起點相減，即使 $|\vec v_1|=|\vec v_2|$（速率相等），$\vec v_2-\vec v_1$ 也\textbf{不是零向量}，因為兩者方向不同。這個差向量 $\Delta\vec v$ 指向圓心。

\textbf{注意：}「加速度」與「變快」並不等同。變快、變慢、轉向三種情形都有加速度；圓周運動是純粹的「轉向」型加速度。可類比為開車以固定速率轉彎：碼錶顯示不變，但乘客仍感覺被推向車門，這即為加速度存在的證據。
\end{補充框}

如下圖：把圓周上相鄰兩點的速度 $\vec v_1$、$\vec v_2$ 平移到一起，畫出 $\Delta\vec v=\vec v_2-\vec v_1$。當兩點靠得很近時，$\Delta\vec v$ 指向圓心。因此這個加速度稱為\textbf{向心加速度（centripetal acceleration）}，centripetal 源自拉丁文「尋求中心」之意。

\fig{si5-centripetal}{0.92}{圖 5-1：左——圓周上速度沿切線方向，速率不變但方向持續改變；右——把相鄰兩點的速度平移相減，差向量 $\Delta\vec v$ 指向圓心，故加速度指向圓心。}

\subsection*{\sffamily B.\ 描述圓周運動的量：$T$、$f$、$\omega$}

\begin{center}
\renewcommand{\arraystretch}{1.3}
\begin{tabular}{p{0.30\linewidth} p{0.10\linewidth} p{0.34\linewidth} p{0.14\linewidth}}
\hline
\textbf{量} & \textbf{符號} & \textbf{定義} & \textbf{單位} \\
\hline
週期（period） & $T$ & 繞一圈所需時間 & s \\
頻率（frequency） & $f$ & 每秒繞幾圈 & Hz \\
角速度（angular velocity） & $\omega$ & 每秒掃過的角度（弧度） & rad/s \\
\hline
\end{tabular}
\end{center}

每繞一圈走過角度 $2\pi$（弧度）、花時間 $T$，因此三者的關係為
$$\boxed{\;f=\frac{1}{T},\qquad \omega=\frac{2\pi}{T}=2\pi f\;}$$

\begin{補充框}{角速度 $\omega$ 是什麼？為什麼用弧度？}
$\omega$ 量的是「轉得多快」（每秒轉幾弧度），與量「跑得多快」的線速率 $v$ 不同，兩者相差一個半徑：$v=\omega r$。

使用\textbf{弧度}而非角度，是為了讓 $v=\omega r$ 這類關係式最簡潔（弧長＝半徑 $\times$ 弧度角）。另須注意 $\omega=2\pi f$，與頻率 $f$ 相差一個 $2\pi$：$f$ 是每秒幾圈，$\omega$ 是每秒幾弧度。
\end{補充框}

\subsection*{\sffamily C.\ 線速率與向心加速度的公式}

繞一圈走過的弧長為圓周長 $2\pi r$，花時間 $T$，因此線速率
$$v=\frac{2\pi r}{T}=\omega r.$$

向心加速度的大小可由「速度三角形」推得。當物體掃過小角度 $\Delta\theta$ 時，速度方向也恰好轉了 $\Delta\theta$。速度大小皆為 $v$，故 $|\Delta\vec v|\approx v\,\Delta\theta$（弧長公式，半徑換成 $v$）。除以 $\Delta t$：
$$a=\frac{|\Delta\vec v|}{\Delta t}=v\,\frac{\Delta\theta}{\Delta t}=v\,\omega.$$
再用 $v=\omega r$ 消去 $\omega$ 或 $v$，得到三個等價寫法：
$$\boxed{\;a=\frac{v^2}{r}=\omega^2 r=v\omega\;}$$

\subsection*{\sffamily D.\ 向心力}

既然有指向圓心的加速度，依牛頓第二定律 $\vec F_{\text{net}}=m\vec a$，必有一個\textbf{指向圓心的合力}維持圓周運動：
$$\boxed{\;F=\frac{mv^2}{r}=m\omega^2 r\;}$$
這個合力稱為\textbf{向心力（centripetal force）}。

\begin{補充框}{向心力不是一種新的力}
\textbf{向心力不是一種新的力，而是「剛好指向圓心的那個合力」的角色名稱。}真正提供它的，永遠是某個已知的真實力——繩的張力、地球的重力、輪胎的靜摩擦力等。向心力只是這些力在圓周運動中所扮演的\textbf{角色}，並非另一個獨立的來源。

牛頓第二定律只指出「合力要等於多少、指向何方」，並未指出這個力\textbf{從何而來}。不同情境的來源各異：

\begin{center}
\renewcommand{\arraystretch}{1.25}
\begin{tabular}{p{0.42\linewidth} p{0.42\linewidth}}
\hline
\textbf{情境} & \textbf{由何者提供向心力} \\
\hline
用繩甩石頭 & 繩的張力 \\
衛星／月球繞地球 & 地球的萬有引力 \\
汽車在水平彎道轉彎 & 輪胎與路面的靜摩擦力 \\
碗裡繞圈的珠子 & 碗面正向力的水平分量 \\
\hline
\end{tabular}
\end{center}

因此畫自由體圖時：\textbf{只畫真實的力（張力、重力、摩擦力、正向力），再把它們的合力設為 $\dfrac{mv^2}{r}$ 指向圓心；不可在真實力之外另畫一支「向心力」箭頭}，否則等於把同一個力計算兩次。

\textbf{注意：}並不存在一個「離心力」把人往外拉。坐車轉彎時感覺被甩向車門，是因為身體依慣性想沿直線前進，而車門把身體拉回圓形軌道；所感覺到的「往外」是慣性的表現，不是真有一支力把人往外推。向心加速度也不會使速率變快，它只負責改變方向。
\end{補充框}

\section*{\sffamily 5-2\quad 簡諧運動}

\subsection*{\sffamily A.\ 回復力}

將彈簧水平放置、一端固定、另一端接一物體，在光滑桌面上拉開後放手，物體會來回振盪。觀察它所受的彈力：物體往右偏（$x>0$）時，彈簧把它往左拉；往左偏（$x<0$）時，彈簧把它往右推。也就是說，力\textbf{永遠指向平衡位置}，且偏離越多力越大。這種力稱為\textbf{回復力（restoring force）}。對理想彈簧，回復力遵守虎克定律：
$$\boxed{\;F=-kx\;}$$
負號表示力的方向與位移 $x$ \textbf{相反}（指回平衡點），$k$ 為彈簧的\textbf{勁度係數（spring constant）}。

\fig{si5-spring}{0.66}{圖 5-2：水平彈簧。上——壓縮（$x<0$）、平衡（$x=0$）、拉伸（$x>0$）三種情形下，回復力 $F=-kx$ 皆指回平衡位置；下——位移隨時間呈餘弦曲線，週期 $T=2\pi\sqrt{m/k}$。}

\textbf{簡諧運動}的定義為：凡是\textbf{回復力與位移成正比、方向相反}（$F=-kx$）的運動。它是最基本的週期運動模型，自然界許多微小振動都近似於它。

\subsection*{\sffamily B.\ 位移、速度、加速度與時間的關係}

把 $F=-kx$ 代入牛頓第二定律 $F=ma$：
$$ma=-kx \;\Rightarrow\; a=-\frac{k}{m}\,x.$$
加速度與位移成正比、方向相反。若放手時物體在最大位移 $A$ 處，其位移隨時間為
$$\boxed{\;x(t)=A\cos(\omega t)\;},\qquad \omega=\sqrt{\frac{k}{m}}.$$
\begin{itemize}\setlength{\itemsep}{1pt}
\item $A$：\textbf{振幅（amplitude）}，離開平衡點的最大距離。
\item $\omega$：\textbf{角頻率（angular frequency）}，決定振盪快慢。
\item $\omega t$（或 $\omega t+\phi$）：\textbf{相位（phase）}，$\phi$ 為初相位，由起始狀態決定。
\end{itemize}

由 $x=A\cos\omega t$ 逐次求變化率（速度是位移的變化率，加速度是速度的變化率）：
$$v(t)=-A\omega\sin(\omega t),\qquad a(t)=-A\omega^2\cos(\omega t)=-\omega^2 x.$$
最後一式 $a=-\omega^2 x$ 與前面 $a=-\dfrac{k}{m}x$ 對照，即得 $\omega^2=\dfrac{k}{m}$。

\begin{補充框}{為什麼解一定是正弦／餘弦？$\omega=\sqrt{k/m}$ 從何而來？}
式子 $a=-\dfrac{k}{m}x$ 表示：這個函數的二次變化率，等於自己乘上一個負數 $-\dfrac{k}{m}$。在基本函數中，只有\textbf{正弦與餘弦}具有這個性質——餘弦取兩次變化率後變回自身的負倍數，正弦亦然。其他函數（如三角形波、拋物線）代入後都無法滿足這個關係。

設 $x=A\cos\omega t$，其二次變化率為 $-A\omega^2\cos\omega t=-\omega^2 x$。與 $a=-\dfrac{k}{m}x$ 比對，立即得到
$$\omega^2=\frac{k}{m}\;\Rightarrow\;\omega=\sqrt{\frac{k}{m}}.$$
因此 $\omega$ 並非外加，而是由方程式形式決定：要讓 $\cos\omega t$ 的二次變化率恰為 $-\dfrac{k}{m}$ 倍，$\omega$ 必須等於 $\sqrt{k/m}$。

\textbf{注意：}餘弦與正弦並非兩種不同的運動，兩者只差一個起始相位（$\sin\omega t=\cos(\omega t-\tfrac{\pi}{2})$）。採用 cos 或 sin，取決於把 $t=0$ 設在最大位移處還是平衡點。
\end{補充框}

由 $\omega=\sqrt{k/m}$ 與 $\omega=\dfrac{2\pi}{T}$，得彈簧簡諧的週期
$$\boxed{\;T=2\pi\sqrt{\frac{m}{k}}\;}$$

兩個結論值得留意：
\begin{itemize}\setlength{\itemsep}{1pt}
\item 週期 $T$ \textbf{只與質量 $m$、勁度 $k$ 有關，與振幅 $A$ 無關}。拉開的幅度不影響來回一趟的時間，此性質稱為\textbf{等時性（isochronism）}。
\item 質量越大、$T$ 越長（振盪越慢）；彈簧越硬（$k$ 大）、$T$ 越短（振盪越快）。
\end{itemize}

\textbf{注意：}振幅不影響週期。拉得越開，路程越長，但速度也越快，兩者恰好抵消。此外，$x=\pm A$（最大位移處）時速度為零、加速度最大（$|a|=\omega^2 A$）；$x=0$（平衡位置）時加速度為零、速度最大（$|v|=\omega A$）。「速度最大」與「加速度最大」發生在不同位置。簡諧的力隨位置變化、不是定值，因此不能套用等加速度三公式。

\subsection*{\sffamily C.\ 簡諧運動的能量}

簡諧系統在振盪時，能量在\textbf{動能（kinetic energy）}與\textbf{位能（potential energy）}之間來回搬運，但\textbf{總力學能不變}。以水平彈簧為例，位能存在被壓縮或拉伸的彈簧裡：
$$U=\tfrac12 kx^2,\qquad K=\tfrac12 mv^2.$$
兩者相加即為總力學能 $E=K+U$。桌面光滑、沒有摩擦做功，因此 $E$ 守恆。用兩個特殊位置就能把 $E$ 定出來：

\begin{itemize}\setlength{\itemsep}{1pt}
\item \textbf{端點}（$x=\pm A$）：此刻速度 $v=0$，能量全是位能，$E=\tfrac12 kA^2$。
\item \textbf{平衡點}（$x=0$）：此刻位能為零、速率最大 $v_{\max}=A\omega$，能量全是動能，$E=\tfrac12 m v_{\max}^2$。
\end{itemize}

兩處的 $E$ 必須相等，於是
$$\boxed{\;E=\tfrac12 kA^2=\tfrac12 m v_{\max}^2=\text{定值}\;}$$
（代入 $v_{\max}=A\omega$ 與 $\omega^2=k/m$，兩式確實相等。）

在任意位置 $x$，由 $E=\tfrac12 kx^2+\tfrac12 mv^2=\tfrac12 kA^2$ 解出速率：
$$\boxed{\;v=\omega\sqrt{A^2-x^2}\;}$$
此式描述整個振盪的快慢分佈：離平衡點越近（$x$ 越小）速率越大，到端點（$x=\pm A$）剛好停住。動能 $K$ 在平衡點最大、端點為零；位能 $U$ 相反，在端點最大、平衡點為零；兩者像接力一樣此消彼長，總和恆為 $\tfrac12 kA^2$。

\begin{補充框}{總能與振幅的關係：$E\propto A^2$}
總力學能 $E=\tfrac12 kA^2$ 與振幅的\textbf{平方}成正比，而非與振幅本身成正比。因此\textbf{振幅加倍時，總能變為 4 倍}（$2^2=4$），不是 2 倍；振幅變 3 倍，總能變 9 倍。

這一點與「週期和振幅無關」並不衝突：拉得越開，振盪一趟的時間相同（等時性），但系統儲存的能量更多、各處速率也更大。週期管的是「快慢的節奏」，總能管的是「振盪的規模」，兩者各自獨立。
\end{補充框}

\textbf{注意：}一個完整振盪裡，動能與位能各自在一個週期內達到最大值兩次（左右兩端點、上下通過平衡點各一次），因此\textbf{動能與位能的變化週期是位移週期的一半}。此外，處理「中途某位置的速率」時，用能量守恆（$\tfrac12 kA^2=\tfrac12 kx^2+\tfrac12 mv^2$）比代入三角函數更直接。

\section*{\sffamily 5-3\quad 簡諧與圓周的關係}

讓一個物體在半徑 $A$ 的圓上做等速圓周運動，角速度為 $\omega$。把它在一條直徑上的投影取出來——\textbf{這個投影的來回運動，即為簡諧運動}。

\fig{si5-projection}{0.92}{圖 5-3：左——半徑 $A$、角速度 $\omega$ 的等速圓周，其在 $x$ 軸上的投影為 $x=A\cos\omega t$；右——投影隨時間描出餘弦曲線，一個來回對應圓周轉一圈，週期同為 $T=2\pi/\omega$。}

設圓周上的點起始於 $x$ 軸正向，轉到角度 $\omega t$ 時，它在 $x$ 軸上的投影為
$$x(t)=A\cos(\omega t).$$
這正是 5-2 的簡諧位移式。因此上一節的解有了明確的幾何來歷。

\begin{補充框}{圓周與簡諧的對應（不是巧合）}
這個對應並非巧合：等速圓周運動的 $x$ 分量本來就滿足 $a_x=-\omega^2 x$（向心加速度的 $x$ 分量），而這正是簡諧的定義式。所以圓周的投影必然是簡諧。

把圓周的速度、加速度向量投影到 $x$ 軸，正好得到簡諧的 $v=-A\omega\sin\omega t$、$a=-A\omega^2\cos\omega t$。正弦／餘弦、$\omega$ 與相位，都可由圓周運動讀出。兩者的對應如下：

\begin{center}
\renewcommand{\arraystretch}{1.25}
\begin{tabular}{p{0.44\linewidth} p{0.44\linewidth}}
\hline
\textbf{圓周運動} & \textbf{投影（簡諧）} \\
\hline
半徑 $r=A$ & 振幅 $A$ \\
角速度 $\omega$ & 角頻率 $\omega$ \\
線速率 $v=\omega A$ & 平衡點最大速率 $v_{\max}=\omega A$ \\
向心加速度 $a=\omega^2 A$ & 端點最大加速度 $a_{\max}=\omega^2 A$ \\
轉一圈時間 $T=\dfrac{2\pi}{\omega}$ & 振盪週期 $T=\dfrac{2\pi}{\omega}$ \\
\hline
\end{tabular}
\end{center}
\end{補充框}

\section*{\sffamily 5-4\quad 單擺}

\subsection*{\sffamily A.\ 小角度下的單擺也是簡諧運動}

\textbf{單擺（simple pendulum）}：一條長 $L$ 的輕線吊一個質量 $m$ 的小球，拉開小角度後放手來回擺動。擺球受兩個力——重力 $mg$（向下）與線張力 $T$（沿線）。把重力沿「擺線方向」與「切線方向」分解：

\fig{si5-pendulum}{0.56}{圖 5-4：單擺受力分解。重力沿擺線的分量 $mg\cos\theta$ 與張力相關，沿切線的分量 $mg\sin\theta$ 指向平衡位置，即為回復力；小角度時 $\sin\theta\approx\theta$，回復力近似正比於位移，故近似為簡諧。}

\begin{itemize}\setlength{\itemsep}{1pt}
\item 沿擺線的分量 $mg\cos\theta$ 與張力相關（使球維持圓弧軌跡）；
\item 沿切線的分量 $mg\sin\theta$ \textbf{指向平衡位置}，即為回復力。
\end{itemize}

回復力大小為 $mg\sin\theta$。當擺角 $\theta$ 很小（約 $\theta<15^\circ$）時，採用\textbf{小角度近似} $\sin\theta\approx\theta$（$\theta$ 取弧度），且弧長位移 $s=L\theta$，於是
$$F_{\text{回}}\approx -mg\theta=-mg\cdot\frac{s}{L}=-\left(\frac{mg}{L}\right)s.$$
這是 $F=-ks$ 的形式，等效勁度 $k=\dfrac{mg}{L}$。代入 $T=2\pi\sqrt{m/k}$：
$$\boxed{\;T=2\pi\sqrt{\frac{L}{g}}\;}$$

由此式可得三個結論：
\begin{itemize}\setlength{\itemsep}{1pt}
\item 週期只與擺長 $L$、重力加速度 $g$ 有關。
\item 週期\textbf{與擺球質量 $m$ 無關}（$m$ 在推導中被約掉）。
\item 週期\textbf{與振幅（擺角）無關}（小角度下），即等時性。
\end{itemize}

\textbf{注意：}$T=2\pi\sqrt{L/g}$ 只在小角度成立。擺角較大時 $\sin\theta\neq\theta$，回復力不再正比於位移，運動不再是嚴格的簡諧，週期會略微變長。此外，「擺長 $L$」是從懸掛點量到\textbf{擺球質心}，若球體較大，須加上其半徑。

\begin{延伸框}{單擺週期為何與振幅無關，且只在小角度成立}
這兩個敘述看似矛盾，實際上是同一個近似的兩面。

\textbf{為什麼振幅大、路程長，週期卻相同？}關鍵在於擺得越開，速度也越快，兩個效應恰好抵消。在簡諧運動中，回復力（即加速度）與位移成正比：拉得越開，回復力越大、加速度越大，球被拉回的速度也越大。路程變長，但速度也變快，一來一回的時間恰好不變。「力正比於位移」是等時性的數學根源——只要力嚴格正比於位移，週期就一定與振幅無關，這對任何簡諧系統（含彈簧）都成立。

\textbf{但單擺的力並非嚴格正比於位移。}單擺真正的回復力是 $mg\sin\theta$，正比的是 $\sin\theta$ 而\textbf{不是} $\theta$（位移 $s=L\theta$）。只有當角度小、$\sin\theta\approx\theta$ 時，回復力才近似正比於位移，才近似為簡諧、才有等時性。因此「振幅無關」是小角度近似的結果：一旦角度大到 $\sin\theta$ 明顯小於 $\theta$，回復力就比簡諧所需的小，週期隨之變長。

精確的單擺週期可展開為
$$T=2\pi\sqrt{\frac{L}{g}}\left(1+\frac{1}{16}\theta_0^2+\cdots\right),$$
其中 $\theta_0$ 為最大擺角（弧度）。$\theta_0$ 很小時修正項 $\dfrac{1}{16}\theta_0^2$ 可忽略，週期幾乎只由 $L$、$g$ 決定（等時性）；$\theta_0$ 變大則修正項變大，週期確實隨振幅增加。例如擺角 $23^\circ$ 時，週期約比小角度公式長 $1\%$。因此等時性嚴格說是\textbf{近似}而非定律，課本「限小角度」即在劃定這個近似的有效範圍。

\textbf{彈簧與單擺的差別}：彈簧的回復力 $F=-kx$ 嚴格正比於位移，等時性對任何振幅都成立；單擺的 $F=mg\sin\theta$ 只是近似正比，僅小角度成立。此為兩者的本質差別。單擺本身是已被完整理解的古典系統，大角度的精確週期可用橢圓積分嚴格表示，上面的級數即為其展開，沒有未解的開放問題。
\end{延伸框}

\section*{\sffamily 5-5\quad 本章重點整理}

\begin{補充框}{一頁複習（公式與關鍵結論）}
\textbf{等速圓周運動}
\begin{itemize}\setlength{\itemsep}{1pt}
\item 速率不變、方向改變，故有\textbf{向心加速度}（指向圓心）。
\item $f=\dfrac1T$，$\omega=\dfrac{2\pi}{T}=2\pi f$，$v=\omega r$。
\item 向心加速度 $a=\dfrac{v^2}{r}=\omega^2 r=v\omega$；向心力 $F=\dfrac{mv^2}{r}=m\omega^2 r$。
\item 向心力\textbf{不是新的力}，由張力／重力／摩擦力等提供；自由體圖不另畫向心力，亦無離心力。
\end{itemize}

\textbf{簡諧運動（SHM）}
\begin{itemize}\setlength{\itemsep}{1pt}
\item 定義：回復力 $F=-kx$（與位移成正比、方向相反）。
\item $x=A\cos\omega t$，$v=-A\omega\sin\omega t$，$a=-\omega^2 x$；$\omega=\sqrt{k/m}$。
\item 端點（$x=\pm A$）：$v=0$、$|a|$ 最大；平衡點（$x=0$）：$|v|=A\omega$ 最大、$a=0$。
\item 彈簧週期 $T=2\pi\sqrt{m/k}$，\textbf{與振幅無關}。
\item \textbf{能量}：$U=\tfrac12 kx^2$、$K=\tfrac12 mv^2$，動能與位能互換而總力學能守恆 $E=\tfrac12 kA^2=\tfrac12 mv_{\max}^2$（$\propto A^2$）；任意位置 $v=\omega\sqrt{A^2-x^2}$。
\end{itemize}

\textbf{簡諧＝圓周的投影}：半徑 $A$、角速度 $\omega$ 的等速圓周，投影到一直徑上即為簡諧；正弦／餘弦與 $\omega$ 由此而來。

\textbf{單擺（小角度）}
\begin{itemize}\setlength{\itemsep}{1pt}
\item 回復力 $\approx mg\sin\theta\approx -\left(\dfrac{mg}{L}\right)s$；$T=2\pi\sqrt{L/g}$。
\item 與質量、振幅無關（僅小角度成立）；可用來測量 $g$。
\end{itemize}
\end{補充框}

\end{document}
